\documentclass[12pt,a4paper]{article}

\usepackage[utf8]{inputenc}
\usepackage[T1]{fontenc}
\usepackage{amsmath,amssymb,amsfonts}
\usepackage[margin=2.5cm]{geometry}
\usepackage{hyperref}
\usepackage{xcolor}
\usepackage{booktabs}
\usepackage[shortlabels]{enumitem}

\newcommand{\Tr}{\operatorname{Tr}}
\newcommand{\tr}{\operatorname{tr}}
\newcommand{\Lam}{\Lambda}
\newcommand{\half}{\tfrac{1}{2}}

\usepackage{amsthm}
\newtheorem{theorem}{Theorem}[section]
\newtheorem{corollary}[theorem]{Corollary}

\definecolor{proven}{RGB}{0,100,0}
\definecolor{near}{RGB}{0,0,160}

\hypersetup{
  pdftitle={Chirality Theorem for Quartic Weyl Structure in the Spectral Action},
  pdfauthor={David Alfyorov},
  colorlinks=true, linkcolor=blue, citecolor=blue, urlcolor=blue
}

\title{\textbf{Chirality Theorem for Quartic Weyl Structure\\[4pt]
in the Spectral Action}}
\author{David Alfyorov}
\date{March 2026}

\begin{document}
\maketitle
\begin{center}
and workflow support.}
\end{center}

\begin{abstract}
We prove that the traced Seeley--DeWitt coefficient $\tr(a_8)$ for the Dirac
operator on a Ricci-flat four-manifold has the form $c(p^2 + q^2)$ with
\emph{zero} cross-term~$pq$, where $p = |C^+|^2$ and $q = |C^-|^2$ are the
norms of the self-dual and anti-self-dual Weyl tensors.  The proof rests on a
chirality theorem: the generators $\sigma^{rs} = \frac{i}{4}[\gamma^r,\gamma^s]$
commute with $\gamma_5$ in $d = 4$, rendering the spin connection curvature
$\Omega_{\mu\nu}$ block-diagonal in the chiral basis.  Consequently the heat
kernel $e^{-tD^2}$ splits into left- and right-handed sectors, each coupled to
only one half of the Weyl tensor (crossed chirality assignment).  This result
applies to the full Standard Model Dirac operator $D$ on the product geometry
$M \times F$ and holds at all orders in the Seeley--DeWitt expansion.
The spectral action at mass dimension~8 therefore generates a single effective
quartic Weyl structure, reducing the three-loop finiteness problem from an
overdetermined system to a single ratio condition.
\end{abstract}

\tableofcontents
\newpage

%===========================================================================
\section{Introduction}
%===========================================================================

The spectral action $S = \Tr(f(D^2/\Lam^2))$ generates, via the Seeley--DeWitt
expansion, curvature invariants at each mass dimension.  At dimension~8 (the
$a_8$ coefficient, equivalently $a_4$ in Avramidi convention), the on-shell
quartic Weyl invariants in $d = 4$ form a two-dimensional space spanned by
$(C^2)^2$ and $({}^*\!CC)^2$.  The FUND program (Sec.~\ref{sec:fund}) established
that the Molien series gives $P(4) = 3$ parity-even quartic Weyl invariants,
reduced to~2 by the Cayley--Hamilton identity.

If the spectral action generates both invariants independently, the single
spectral parameter $\delta f_8$ cannot absorb a generic three-loop counterterm
with two independent coefficients.  The question is whether the spectral action's
specific structure constrains the quartic Weyl ratio.

\section{The Chirality Theorem}
\label{sec:chirality}

\begin{theorem}[Chirality of the spin connection]
In $d = 4$ Euclidean (or Lorentzian) dimensions, the spin connection generators
$\sigma^{rs} = \frac{i}{4}[\gamma^r, \gamma^s]$ commute with the chirality
operator~$\gamma_5$:
\begin{equation}\label{eq:sigma-g5}
  [\sigma^{rs},\, \gamma_5] = 0 \quad \forall\; r, s \in \{0,1,2,3\}.
\end{equation}
\end{theorem}

\begin{proof}
The chirality operator $\gamma_5 = \gamma^0\gamma^1\gamma^2\gamma^3$ satisfies
$\{\gamma_5, \gamma^\mu\} = 0$ for each $\mu$.  Therefore
$\gamma_5\, \gamma^r\gamma^s = (-\gamma^r\gamma_5)\gamma^s
= \gamma^r\gamma^s\,\gamma_5$, giving
$[\gamma_5, \gamma^r\gamma^s] = 0$.  Since
$\sigma^{rs} \propto [\gamma^r, \gamma^s] = \gamma^r\gamma^s - \gamma^s\gamma^r$,
and both products commute with $\gamma_5$, the result follows.
\end{proof}

\begin{corollary}[Block-diagonal spin connection curvature]
The spin connection curvature
$\Omega_{\mu\nu} = \frac{1}{4}R_{\mu\nu\rho\sigma}\sigma^{\rho\sigma}$
commutes with $\gamma_5$ and is block-diagonal in the chiral basis
$\gamma_5 = \mathrm{diag}(I_2, -I_2)$:
\begin{equation}
  \Omega_{\mu\nu} = \begin{pmatrix} \Omega^L_{\mu\nu} & 0 \\ 0 & \Omega^R_{\mu\nu} \end{pmatrix}.
\end{equation}
\end{corollary}

\begin{corollary}[Crossed chirality assignment]
\label{cor:crossed}
The left-handed block $\Omega^L$ couples exclusively to the anti-self-dual
Weyl tensor $C^-$ (parameter~$q$), and the right-handed block $\Omega^R$
couples exclusively to the self-dual Weyl tensor $C^+$ (parameter~$p$):
\begin{equation}
  \tr_L(\Omega^2) = -\half\, q, \qquad
  \tr_R(\Omega^2) = -\half\, p.
\end{equation}
This follows from the 't~Hooft symbol assignment:
$\eta^i_{ab}\sigma^{ab}$ lives in~$P_R$ and
$\bar\eta^i_{ab}\sigma^{ab}$ lives in~$P_L$.
\end{corollary}

\section{Heat Kernel Decomposition}
\label{sec:heat}

On a Ricci-flat background ($R_{\mu\nu} = 0$), the Lichnerowicz--Weitzenb\"ock
formula gives $D^2 = -\nabla^2_{\mathrm{spin}}$ with endomorphism $E = -R/4 = 0$.
Since the spin connection $\nabla_\mu = \partial_\mu + \Gamma_\mu$ with
$\Gamma_\mu = \frac{1}{4}\omega_\mu^{ab}\gamma_a\gamma_b$ commutes
with~$\gamma_5$ (Theorem~1), the operator~$D^2$ commutes with~$\gamma_5$:
\begin{equation}
  [D^2,\, \gamma_5] = 0 \quad \text{on Ricci-flat } M^4.
\end{equation}
Therefore $e^{-tD^2}$ is block-diagonal in the chiral basis, and the traced
heat kernel coefficients decompose:
\begin{equation}
  \tr(a_n) = \tr_L(a_n^L) + \tr_R(a_n^R)
\end{equation}
where $a_n^L$ depends only on~$q$ (anti-self-dual Weyl) and $a_n^R$ depends
only on~$p$ (self-dual Weyl).

\begin{theorem}[Quartic Weyl ratio]
At the quartic level ($n = 4$ in Avramidi convention, corresponding to mass
dimension~8):
\begin{equation}\label{eq:ratio}
  \tr(a_8) = c\,(p^2 + q^2) = \frac{c}{2}\bigl[(C^2)^2 + ({}^*\!CC)^2\bigr],
\end{equation}
with zero $pq$ cross-term.  The ratio
$(C^2)^2 : ({}^*\!CC)^2 = 1{:}1$ is exact.
\end{theorem}

\begin{proof}
By the chiral decomposition: $\tr(a_8) = f_4(p) + f_4(q)$ where $f_4$ is a
degree-2 polynomial (quartic in Weyl = degree~2 in $p$ or~$q$).  By parity
($p \leftrightarrow q$): $f_4(p) = f_4(q)$, so $f_4(x) = c\,x^2$ for some
constant~$c$.  Therefore $\tr(a_8) = c(p^2 + q^2)$.
\end{proof}

\section{Extension to the Full Standard Model}
\label{sec:sm}

The Standard Model Dirac operator on the almost-commutative geometry
$M \times F$ takes the form:
\begin{equation}
  D = D_M \otimes \mathbf{1}_F + \gamma_5 \otimes D_F.
\end{equation}
On a pure-gravity background (no gauge field fluctuations):
\begin{equation}
  D^2 = D_M^2 \otimes \mathbf{1}_F + \{D_M, \gamma_5\} \otimes D_F
  + \mathbf{1} \otimes D_F^2.
\end{equation}
In $d = 4$: $\{D_M, \gamma_5\} = i\{\gamma^\mu, \gamma_5\}\nabla_\mu = 0$
(since $\{\gamma^\mu, \gamma_5\} = 0$).  Therefore:
\begin{equation}
  [D^2,\, \gamma_5 \otimes \mathbf{1}_F] = [D_M^2, \gamma_5] \otimes \mathbf{1}_F = 0.
\end{equation}
The chirality theorem extends to the full SM spectral triple.  The spectral action
$\Tr(f(D^2/\Lam^2))$ generates quartic Weyl invariants with ratio~$1{:}1$ at
all orders in the Seeley--DeWitt expansion.

\section{Quartic $\Omega$-Structure Decomposition}
\label{sec:omega}

Four quartic $\Omega$-structures appear in the heat kernel:

\begin{center}
\begin{tabular}{lccc}
\toprule
Structure & $(p^2{+}q^2)$ coeff & $pq$ coeff & Block-diagonal? \\
\midrule
$S_1 = \tr(\Omega^4_{\mathrm{chain}})$ & $1/16$ & $0$ & Yes \\
$S_2 = \tr(\Omega_{\mathrm{sq}}^2)$ & $1/8$ & $0$ & Yes \\
$S_3 = [\tr(\Omega^2)]^2$ & $1/4$ & $1/2$ & \textbf{No} \\
$S_4 = \text{double trace}$ & $1/8$ & $0$ & Yes \\
\bottomrule
\end{tabular}
\end{center}

Only $S_3$ has a nonzero $pq$ coefficient.  However, the chirality theorem
guarantees that the \emph{total} $\tr(a_8)$ has zero~$pq$.  This implies
a non-trivial cancellation: the $pq$ from~$S_3$ (with its universal Avramidi
coefficient) must cancel against $pq$ contributions from geometry-sector
terms (pure Riemann$^4 \times \mathbf{1}$ in fiber space).  The cancellation
is forced by the chiral block structure of each sector's heat kernel.

\section{Implications for the Three-Loop Problem}
\label{sec:fund}

The spectral action at mass dimension~8 generates:
\begin{equation}
  S_8 = f_8 \cdot c \int d^4x\,\sqrt{g}\; (p^2 + q^2)
  = \frac{f_8 c}{2} \int d^4x\,\sqrt{g}\;
  \bigl[(C^2)^2 + ({}^*\!CC)^2\bigr].
\end{equation}
This is a \emph{single} effective structure controlled by one parameter~$f_8$.
The three-loop counterterm absorption requires:
\begin{equation}
  \delta f_8 \cdot c = \alpha_{\mathrm{ct}} \quad \text{(coefficient of } p^2{+}q^2)
\end{equation}
which is one equation in one unknown---always solvable---provided the
three-loop counterterm also has the form $\alpha_{\mathrm{ct}}(p^2 + q^2)$
with zero~$pq$.

\subsection{Status of the counterterm ratio}

The chirality theorem proves $pq = 0$ for:
\begin{itemize}
  \item The spectral action heat kernel $\tr(a_n)$ at all orders (proven).
  \item The one-loop counterterm (which \emph{is} $\tr(a_8)$, proven).
\end{itemize}
For multi-loop counterterms ($L \geq 2$), the argument requires that the
spectral action class is preserved under renormalization.  van~Nuland and
van~Suijlekom (arXiv:2107.08485) proved this for Yang--Mills on flat space.
Extension to the gravitational sector is supported by:
\begin{enumerate}[(i)]
  \item All vertices $\delta^n S/\delta g^n$ inherit chirality from~$D^2$.
  \item The kinetic operator $\delta^2 S/\delta h^2$ acts on symmetric
    tensors in the $(3,3)$ representation of $\mathrm{SU}(2)_L \times \mathrm{SU}(2)_R$,
    which admits a chiral decomposition.
  \item The propagator $G = K^{-1}$ inherits the block structure if $K$
    preserves it.
\end{enumerate}
A formal proof for the gravitational sector remains open.

\section{Conclusion}

The chirality theorem reduces the three-loop finiteness problem from a
structurally overdetermined system ($P(4) = 3$ invariants, reduced to~2 by
Cayley--Hamilton, versus~1 spectral parameter) to a single ratio condition
with zero free parameters.  The spectral action side is proven to generate
only the $(p^2 + q^2)$ combination.  Resolution of the three-loop problem
is conditional on the counterterm sharing this structure, which is guaranteed
if the spectral action class is renormalizably closed for gravity.

\end{document}
